% Macros and settings used by both the pdf and the web version.

% Theorem-like environments (numbered within chapters). The dependency graph draws definition,
% lemma, proposition, theorem and corollary (the default of the `thms=` option in web.tex), not
% remark.
\newtheorem{theorem}{Theorem}[chapter]
\newtheorem{proposition}[theorem]{Proposition}
\newtheorem{lemma}[theorem]{Lemma}
\newtheorem{corollary}[theorem]{Corollary}

\theoremstyle{definition}
\newtheorem{definition}[theorem]{Definition}
\newtheorem{remark}[theorem]{Remark}

% Notation.
\newcommand{\QQ}{\mathbb{Q}}
\newcommand{\ZZ}{\mathbb{Z}}
\newcommand{\NN}{\mathbb{N}}

% The name of the C program and library, in typewriter like the name of the executable
% ratpoints-lean, the options and the files; the bookmarks of the PDF take the plain name.
\newcommand{\ratpoints}{\texorpdfstring{\texttt{ratpoints}}{ratpoints}}

% The ratpoints commit the line numbers of the C code refer to: the pin of
% correspondence/ratpoints.toml (scripts/correspondence.py checks that the two agree).
\newcommand{\ratpointscommit}{1a998538dfa835bb1bce6047e75302cc01517443}
% A link to lines #2 to #3 of the C file #1 of ratpoints at that commit. Inside a node of the
% dependency graph that names declarations with \lean, every \csource link must be one of the
% ranges named in the `C:` paragraphs of the docstrings of the model declarations it names;
% prose and remarks may link to any range. scripts/correspondence.py checks this, and every link
% is locked with the text it shows (correspondence/README.md).
\newcommand{\csource}[3]{\href{https://github.com/MichaelStollBayreuth/ratpoints/blob/\ratpointscommit/#1\#L#2-L#3}{\nolinkurl{#1:#2-#3}}}

% Marks a place where the formal statement or the model deviates from the C code or the manual.
\newcommand{\formalnote}[1]{\par\smallskip\noindent\textit{Formalization:} #1\par\smallskip}
